\documentclass[11pt,letterpaper]{article}
\usepackage[T1]{fontenc}
\usepackage[utf8]{inputenc}
\usepackage{lmodern,microtype}
\usepackage[margin=0.78in,headheight=15pt]{geometry}
\usepackage{amsmath,amssymb,booktabs,longtable,tabularx,array,enumitem,xcolor,hyperref,fancyhdr,lastpage,needspace}
\definecolor{EngineBlue}{HTML}{183B56}
\definecolor{EngineTeal}{HTML}{117A83}
\definecolor{EngineGray}{HTML}{596671}
\definecolor{EngineLight}{HTML}{EEF4F7}
\hypersetup{colorlinks=true,urlcolor=EngineTeal,linkcolor=EngineBlue}
\setlength{\parindent}{0pt}
\setlength{\parskip}{5pt}
\emergencystretch=3em
\setlist{nosep,leftmargin=*}
\pagestyle{fancy}\fancyhf{}
\fancyhead[L]{\small\sffamily Game Engine Architecture | Engineering Suite}
\fancyhead[R]{\small\sffamily \DocType}
\fancyfoot[L]{\scriptsize Source-grounded engineering curriculum | v1.0}
\fancyfoot[R]{\scriptsize \thepage/\pageref{LastPage}}
\newcommand{\DocType}{Reference}
\newcommand{\SuiteHeading}[3]{%
\renewcommand{\DocType}{#2}%
{\small\bfseries\sffamily\color{EngineTeal}#2\hfill #3}\par\vspace{6pt}
{\LARGE\bfseries\color{EngineBlue}#1}\par\vspace{4pt}
{\color{EngineTeal}\rule{\linewidth}{1pt}}\par\vspace{7pt}}
\newcommand{\LabEvidence}[1]{\par\textbf{Evidence required:} #1\par}
\newcommand{\WriteLines}[1]{\par\begingroup\setlength{\parskip}{9pt}\count255=0\loop\ifnum\count255<#1\noindent\rule{\linewidth}{0.35pt}\par\advance\count255 by1\repeat\endgroup\smallskip}
\newcommand{\WritingBox}[1]{\noindent\fbox{\parbox[t][#1][t]{0.97\linewidth}{\mbox{}}}\par\smallskip}
\newcommand{\SourceNote}{\footnotesize\textit{Primary reading: Jason Gregory, \emph{Game Engine Architecture}, 3rd ed. (CRC Press, 2018). The narrative, exercises, laboratory designs and scoring instruments in this suite are independently authored teaching materials, not a reproduction of the textbook.}\normalsize\par}
\newcommand{\Criterion}[2]{\textbf{#1.} #2\par}
\newcommand{\FormField}[2]{\Needspace{3\baselineskip}\textbf{#1}\par\textcolor{EngineGray}{\footnotesize #2}\par\WriteLines{2}}
\newcommand{\BlankField}[1]{\Needspace{3\baselineskip}\textbf{#1}\par\WriteLines{2}}
\newcolumntype{Y}{>{\raggedright\arraybackslash}X}

\begin{document}
\SuiteHeading{Job Scheduler Scaling and Correctness Procedure}{Benchmark procedure}{Evidence protocol}
\section{Objective and conditions}
Procedure: benchmark original laboratory implementations under repeatable test conditions. Treat each measurement as a property of the exact workload and environment; any comparison across hosts must state its limitations.

\subsection{Reference}
Use a single-thread deterministic output and checksum for the same data.

\subsection{Workloads}
Sweep worker counts 1,2,4,8 (up to available hardware), job granularity, computational intensity and shared-state contention.

\subsection{Correctness}
Run repeated schedules with randomized task order and verify exactly-once completion, shutdown and exception paths.

\subsection{Performance}
Report speedup, throughput, p95 per-batch latency, worker idle time and oversubscription impact.

\subsection{Limitations}
Identify memory bandwidth saturation and scheduler overhead; avoid bitwise checks for order-sensitive floating-point reduction unless ordered.

\section{Required CSV schema}
\begin{verbatim}
experiment,revision,host,variant,seed,iteration,elapsed_ns,
work_items,peak_bytes,success,notes
\end{verbatim}\section{Report checklist}
\begin{itemize}
\item Input and output artifact hashes
\item At least one cold and warm run where applicable
\item Raw data, median and tail-percentile summary
\item Error bars or limitations for small sample sizes
\item Explicit negative tests and documented release decision
\end{itemize}

\end{document}
